
\documentclass{article}
\usepackage{colortbl}
\usepackage{makecell}
\usepackage{multirow}
\usepackage{supertabular}

\begin{document}

\newcounter{utterance}

\centering \large Interaction Transcript for game `dond', experiment `coop\_en', episode 16 with reflexion{-}Qwen3.5{-}4B.
\vspace{24pt}

{ \footnotesize  \setcounter{utterance}{1}
\setlength{\tabcolsep}{0pt}
\begin{supertabular}{c@{$\;$}|p{.15\linewidth}@{}p{.15\linewidth}p{.15\linewidth}p{.15\linewidth}p{.15\linewidth}p{.15\linewidth}}
    \# & $\;$A & \multicolumn{4}{c}{Game Master} & $\;\:$B\\
    \hline

    \theutterance \stepcounter{utterance}  
    & & \multicolumn{4}{p{0.6\linewidth}}{
        \cellcolor[rgb]{0.9,0.9,0.9}{
            \makecell[{{p{\linewidth}}}]{
                \texttt{\tiny{[P1$\langle$GM]}}
                \texttt{You have attempted this same task before and did not succeed. These are the plans you wrote after those attempts.} \\
\texttt{Trial \#1: New plan: I failed because I misunderstood the game mechanics, specifically the distinction between communicating a division plan to the opponent versus submitting a secret proposal. I repeatedly sent "magnet, carrot" as if it were a proposal, which the system interpreted as a secret proposal, causing the game to lock or time out before a valid agreement could be reached. I should have first communicated a clear, non{-}secret proposal to the opponent (e.g., "I take the magnet and carrot; you take the leaves and wire") to establish the agreement, and only then submitted the formal secret proposal in the exact required format: "{[}Proposal: 1 magnet, 1 carrot{]}". I must ensure that any message sent to the opponent does not use square brackets, reserving that format strictly for the final submission step.} \\
\\ 
\texttt{Trial \#2: New plan: I failed because I submitted a secret proposal before the other player had submitted theirs, violating the turn{-}based requirement for complementary proposals. The game requires both players to submit proposals before a valid agreement can be reached. Next time, I will first send a clear, non{-}secret message to the opponent outlining the division (e.g., "I take the magnet and carrot; you take the leaves and wire") to propose the agreement. Only after receiving the opponent's response or confirmation will I submit my formal secret proposal using the exact format "{[}Proposal: 1 magnet, 1 carrot{]}" to lock in the deal. I must ensure I do not submit a proposal until the opponent has also submitted one.} \\
\\ 
\\ 
\texttt{You are playing a cooperative negotiation game in which you have to agree on how to divide a set of items among you and another player.} \\
\\ 
\texttt{Rules:} \\
\texttt{(a) You and the other player are given a set of items. Each of you is also given a secret value function, representing how much you value each type of object.} \\
\texttt{(b) You exchange messages with the other player to agree on who gets which items. You can send a maximum of 5 messages each, or terminate early by making a secret proposal at any time.} \\
\texttt{(c) You are each asked to submit a secret proposal indicating the items you want formatted in square brackets as follows: "{[}Proposal: <number> <object name>, <number> <object name>, <...>{]}"} \\
\texttt{(d) If your proposals are complementary, i.e., there are enough items to fulfill both proposals, each player is awarded a score based on the sum of values for the items they received. Otherwise, both of you get zero points.} \\
\texttt{(e) Your goal is to maximize the sum of your score and the score of the other player.} \\
\\ 
\texttt{Let us start.} \\
\\ 
\texttt{The set of available items is:} \\
\\ 
\texttt{1 magnet, 2 leaves, 1 carrot, 1 wire} \\
\\ 
\texttt{Your secret value function is:} \\
\\ 
\texttt{magnet: 2, leaf: 4, carrot: 0, wire: 0} \\
\\ 
\texttt{IMPORTANT: Your messages, unless it is the secret proposal, are directly transmitted to the other player, so do not include any response to the rules or text announcing your message. To make a secret proposal, use the indicated format. Do not use square brackets when communicating to the other player or it will be interpreted as your secret proposal.} \\
            }
        }
    }
    & \\ \\

    \theutterance \stepcounter{utterance}  
    & \multicolumn{4}{p{0.6\linewidth}}{
        \cellcolor[rgb]{0.9,0.9,0.9}{
            \makecell[{{p{\linewidth}}}]{
                \texttt{\tiny{[P1$\rangle$GM]}}
                \texttt{I take the magnet and carrot; you take the leaves and wire.} \\
            }
        }
    }
    & & \\ \\

    \theutterance \stepcounter{utterance}  
    & & & \multicolumn{2}{p{0.3\linewidth}}{
        \cellcolor[rgb]{0.9,0.9,0.9}{
            \makecell[{{p{\linewidth}}}]{
                \texttt{\tiny{[GM$|$GM]}}
                \texttt{continue} \\
            }
        }
    }
    & & \\ \\

    \theutterance \stepcounter{utterance}  
    & & & \multicolumn{2}{p{0.3\linewidth}}{
        \cellcolor[rgb]{0.9,0.9,0.9}{
            \makecell[{{p{\linewidth}}}]{
                \texttt{\tiny{[GM$|$GM]}}
                \texttt{I take the magnet and carrot; you take the leaves and wire.} \\
            }
        }
    }
    & & \\ \\

    \theutterance \stepcounter{utterance}  
    & & \multicolumn{4}{p{0.6\linewidth}}{
        \cellcolor[rgb]{0.9,0.9,0.9}{
            \makecell[{{p{\linewidth}}}]{
                \texttt{\tiny{[GM$\rangle$P2]}}
                \texttt{You have attempted this same task before and did not succeed. These are the plans you wrote after those attempts.} \\
\texttt{Trial \#1: New plan: I failed because I misunderstood the game's turn structure and the definition of a "secret proposal." I incorrectly assumed that sending a list of items (e.g., "magnet, carrot") without the specific bracketed format `{[}Proposal: ...{]}` would be a valid negotiation move or a hidden proposal. In reality, the rules state that only the bracketed format constitutes a secret proposal, and all other messages are direct communications to the opponent. By sending unformatted text, I likely triggered a state where the opponent's subsequent unformatted message was interpreted as a proposal, or I simply exhausted my turns without ever making a valid, complementary offer. The core mistake was treating any item list as a proposal. The corrected strategy is to strictly adhere to the communication protocol: use only the specific `{[}Proposal: <number> <object>, ...{]}` format to submit an offer. Before submitting, I must calculate a distribution where the sum of my values plus the opponent's values (assuming they value the items I don't have) is maximized, ensuring the items I propose are exactly what the opponent needs to reach a complementary state. I will start by analyzing the item values (magnet: 3, leaf: 1, carrot: 4, wire: 1) and propose a split that gives the opponent the high{-}value items they likely desire (e.g., the carrot) while I take the magnet, ensuring the proposal format is perfect.} \\
\\ 
\texttt{Trial \#2: New plan: I failed because I submitted a proposal that was not complementary to any possible response from the opponent. My proposal took the magnet (value 3) and carrot (value 4), totaling 7. The remaining items were two leaves (value 1 each) and one wire (value 1), totaling 2. Even if the opponent valued the leaves and wire highly, they could not reach a complementary state with my specific set because the remaining items (leaves, wire) do not form a complete set for a second player if I take the magnet and carrot, unless the opponent's proposal exactly matches the leftovers. However, the critical error was likely that I did not verify if the *remaining* items could satisfy the opponent's potential high{-}value needs or if my own proposal left the opponent with nothing of value they needed. Specifically, I should have checked if the opponent could propose a set that, combined with mine, covers all items. Since I took the carrot (my highest value) and magnet (high value), I left the opponent with only leaves and wire. If the opponent values the carrot more than I do (or values the magnet more), my proposal fails. The correct strategy is to propose a split where the items I propose are exactly the complement of the items the opponent *must* take to get a valid score. Given the values (Carrot: 4, Magnet: 3, Leaves: 1, Wire: 1), a successful complementary split requires one player to take the Carrot and the other to take the Magnet (and perhaps some leaves/wire). I should have proposed taking the Carrot and one Leaf (or Wire) to leave the Magnet and the other Leaf/Wire for the opponent, ensuring the sum of values is maximized and the sets are disjoint and cover the necessary items. I will now propose: {[}Proposal: 1 carrot, 1 leaf{]} to leave the magnet and wire for the opponent, or better yet, {[}Proposal: 1 carrot, 1 wire{]} to leave the magnet and leaves. Actually, the most robust complementary split is me taking the Carrot and Magnet (7 points) and the opponent taking the Leaves and Wire (2 points), but only if the opponent values the leaves/wire enough to accept. Wait, the rule says "If your proposals are complementary... each player is awarded a score". This implies I need to propose a set, and the opponent proposes a set, and if Set\_A + Set\_B = All Items, we win. My previous proposal was {[}1 magnet, 1 carrot{]}. For this to be complementary, the opponent must propose {[}2 leaves, 1 wire{]}. If the opponent does that, we win. Why did I lose? Perhaps the opponent did not propose exactly that, or I misjudged the likelihood. The real mistake was assuming a specific split without verifying the complementarity condition explicitly in my head before sending. I will now propose a} \\
\\ 
\\ 
\texttt{You are playing a cooperative negotiation game in which you have to agree on how to divide a set of items among you and another player.} \\
\\ 
\texttt{Rules:} \\
\texttt{(a) You and the other player are given a set of items. Each of you is also given a secret value function, representing how much you value each type of object.} \\
\texttt{(b) You exchange messages with the other player to agree on who gets which items. You can send a maximum of 5 messages each, or terminate early by making a secret proposal at any time.} \\
\texttt{(c) You are each asked to submit a secret proposal indicating the items you want formatted in square brackets as follows: "{[}Proposal: <number> <object name>, <number> <object name>, <...>{]}"} \\
\texttt{(d) If your proposals are complementary, i.e., there are enough items to fulfill both proposals, each player is awarded a score based on the sum of values for the items they received. Otherwise, both of you get zero points.} \\
\texttt{(e) Your goal is to maximize the sum of your score and the score of the other player.} \\
\\ 
\texttt{Let us start.} \\
\\ 
\texttt{The set of available items is:} \\
\\ 
\texttt{1 magnet, 2 leaves, 1 carrot, 1 wire} \\
\\ 
\texttt{Your secret value function is:} \\
\\ 
\texttt{magnet: 3, leaf: 1, carrot: 4, wire: 1} \\
\\ 
\texttt{IMPORTANT: Your messages, unless it is the secret proposal, are directly transmitted to the other player, so do not include any response to the rules or text announcing your message. To make a secret proposal, use the indicated format. Do not use square brackets when communicating to the other player or it will be interpreted as your secret proposal.} \\
\\ 
\\ 
\texttt{I take the magnet and carrot; you take the leaves and wire.} \\
            }
        }
    }
    & \\ \\

    \theutterance \stepcounter{utterance}  
    & & & \multicolumn{4}{p{0.6\linewidth}}{
        \cellcolor[rgb]{0.9,0.9,0.9}{
            \makecell[{{p{\linewidth}}}]{
                \texttt{\tiny{[GM$\langle$P2]}}
                \texttt{{[}Proposal: 1 magnet, 1 carrot{]}} \\
            }
        }
    }
     \\ \\

    \theutterance \stepcounter{utterance}  
    & & & \multicolumn{2}{p{0.3\linewidth}}{
        \cellcolor[rgb]{0.9,0.9,0.9}{
            \makecell[{{p{\linewidth}}}]{
                \texttt{\tiny{[GM$|$GM]}}
                \texttt{proposal} \\
            }
        }
    }
    & & \\ \\

    \theutterance \stepcounter{utterance}  
    & & & \multicolumn{2}{p{0.3\linewidth}}{
        \cellcolor[rgb]{0.9,0.9,0.9}{
            \makecell[{{p{\linewidth}}}]{
                \texttt{\tiny{[GM$|$GM]}}
                [1, 0, 1, 0]
            }
        }
    }
    & & \\ \\

    \theutterance \stepcounter{utterance}  
    & & \multicolumn{4}{p{0.6\linewidth}}{
        \cellcolor[rgb]{0.9,0.9,0.9}{
            \makecell[{{p{\linewidth}}}]{
                \texttt{\tiny{[P1$\langle$GM]}}
                \texttt{The other player has submitted a secret proposal. It is now your turn to submit a proposal of your own.} \\
            }
        }
    }
    & \\ \\

    \theutterance \stepcounter{utterance}  
    & \multicolumn{4}{p{0.6\linewidth}}{
        \cellcolor[rgb]{0.9,0.9,0.9}{
            \makecell[{{p{\linewidth}}}]{
                \texttt{\tiny{[P1$\rangle$GM]}}
                \texttt{{[}Proposal: 1 magnet, 1 carrot{]}} \\
            }
        }
    }
    & & \\ \\

    \theutterance \stepcounter{utterance}  
    & & & \multicolumn{2}{p{0.3\linewidth}}{
        \cellcolor[rgb]{0.9,0.9,0.9}{
            \makecell[{{p{\linewidth}}}]{
                \texttt{\tiny{[GM$|$GM]}}
                \texttt{proposal} \\
            }
        }
    }
    & & \\ \\

    \theutterance \stepcounter{utterance}  
    & & & \multicolumn{2}{p{0.3\linewidth}}{
        \cellcolor[rgb]{0.9,0.9,0.9}{
            \makecell[{{p{\linewidth}}}]{
                \texttt{\tiny{[GM$|$GM]}}
                [1, 0, 1, 0]
            }
        }
    }
    & & \\ \\

    \theutterance \stepcounter{utterance}  
    & & & \multicolumn{2}{p{0.3\linewidth}}{
        \cellcolor[rgb]{0.9,0.9,0.9}{
            \makecell[{{p{\linewidth}}}]{
                \texttt{\tiny{[GM$|$GM]}}
                [[1, 0, 1, 0], [1, 0, 1, 0]]
            }
        }
    }
    & & \\ \\

    \theutterance \stepcounter{utterance}  
    & \multicolumn{4}{p{0.6\linewidth}}{
        \cellcolor[rgb]{0.95,0.95,0.95}{
            \makecell[{{p{\linewidth}}}]{
                \texttt{\tiny{[P2$\rightarrow$P1]}}
                \texttt{Post{-}episode Reflexion (not sent to GM):} \\
\texttt{New plan: I failed because I submitted my secret proposal immediately after the opponent's, without verifying that our proposals were complementary or that the opponent had actually accepted the division. The rules state that if proposals are not complementary, both get zero points. My previous proposal only claimed items I wanted, but I did not confirm if the opponent's proposal left exactly what I needed. Next time, after the opponent submits their proposal, I will first analyze the combined set of items in both proposals to check for complementarity. If they are complementary, I will immediately submit my formal secret proposal in the correct format "{[}Proposal: ...{]}" to lock the deal. If they are not, I will send a direct message to the opponent suggesting a revised division that ensures both parties get their desired items, rather than submitting a conflicting secret proposal. I must strictly avoid submitting a secret proposal until I have confirmed the other player's proposal and verified the agreement is valid.} \\
            }
        }
    }
    & & \\ \\

    \theutterance \stepcounter{utterance}  
    & & & \multicolumn{4}{p{0.6\linewidth}}{
        \cellcolor[rgb]{0.95,0.95,0.95}{
            \makecell[{{p{\linewidth}}}]{
                \texttt{\tiny{[P1$\rightarrow$P2]}}
                \texttt{Post{-}episode Reflexion (not sent to GM):} \\
\texttt{New plan: I failed because I submitted a proposal for only two items (magnet and carrot) without ensuring that the remaining items (leaves and wire) could form a valid, complementary set for the opponent to propose. The rules state that proposals are complementary only if there are enough items to fulfill both; my partial proposal left the opponent with a set that might not be sufficient or desirable enough to trigger a successful match, or I simply did not propose a complete partition of the items. To succeed, I must propose a set of items that, when combined with the opponent's potential proposal, covers the entire set of items exactly once. Given the items (1 magnet, 2 leaves, 1 carrot, 1 wire), I should propose a set that leaves the remaining items as a coherent, complete set for the opponent. For instance, if I propose taking the magnet and one leaf, the opponent can propose taking the carrot, the other leaf, and the wire. Alternatively, I could propose taking the carrot and the magnet, but I must be certain the opponent will propose the exact remaining set (2 leaves, 1 wire). The safest strategy is to propose a split that clearly defines a complete remainder. I will now propose: {[}Proposal: 1 magnet, 1 leaf{]} to leave the carrot, 1 leaf, and 1 wire for the opponent, ensuring the sets are disjoint and cover all items.} \\
            }
        }
    }
     \\ \\

\end{supertabular}
}

\end{document}
